We study the numerical evaluation of contour integral representations for
solutions of the linearized second and fourth Painlev\'e equations.
The construction combines a nonlinear background solution, continuation of
canonical Lax-pair columns, normalization matching, and quadrature over
cycles with decaying branches in distinct sectors. Accuracy is assessed
using a reference fundamental matrix, differential-equation residuals,
and variations of monodromy data computed independently from the spectral
problem. An exact identity describes the propagation of initial-basis
errors. A stagewise a posteriori refinement criterion uses increments in
monodromy variations to allocate numerical accuracy. The contour algorithm
is compared with the Dormand--Prince method on a regular background, in
a region of rapid solution growth, and through hard loss of stability
followed by approximately regular oscillations. Accumulated offsets in
monodromy variations and subsequent drift are measured separately.
The experiments show how drift depends on spectral conditioning,
quadrature accuracy, and error allocation between initial and subsequent
parts of the computation. Five additional initial-data and parameter
cases demonstrate the dependence of comparative performance on the
nonlinear background.
